\subsection{Cross-model contrast does not isolate a topology law}
Verdict \#6 asks whether the fragility framework applies to a second developmental network. We applied the same basin-enumeration and single-bit-exit construction to the four-node gap-gene cascade with an anterior Bcd prepattern. This implemented module has one global attractor: its WT basin fraction is 1.0 and its basin-average one-bit exit fraction is zero. The implemented segment-polarity network has a WT basin fraction of 0.781\% and exit fraction 0.500. The numerical contrast is reproducible, but two different models do not isolate topology from their regulatory rules, external clamps, state definitions or update schedules.

We therefore withdraw the earlier claim that fragility is a property of feedback topology rather than modeling choices. The floral experiments in this paper supply direct counterexamples to that wording: with fixed topology, fate rankings change with starting measure, asynchronous update and recovery horizon. A single-global-attractor system has zero eventual basin-exit fragility by definition under its deterministic rule, not because every cascade or real developmental network is globally robust. Feedback circuits can have several attractors, but these examples do not prove that feedback necessarily causes a given degree of fragility.

The useful transferable object is the specified measurement protocol: identify the transition rule, state measure, perturbation kernel and horizon, then compare robustness under controlled changes. The two-model contrast motivates a topology hypothesis, not a validated causal law. Original numerical records remain in \path{results/grn_fragility_gapgene.json} and \path{experiments/fragility_gapgene.py}; their scope is unchanged.
